\usepackage{tikz}
\usepackage{xcolor}
\usepackage{xparse}
\usepackage{ifthen}
\usepackage{pgfkeys}
\usetikzlibrary{positioning} % Required for 'of' positioning
\usetikzlibrary{arrows.meta} % For arrows

\definecolor{green}{HTML}{6ca877}
\definecolor{purple}{HTML}{805695}
\definecolor{yellow}{HTML}{e89959}
\definecolor{red}{HTML}{e06262}
\definecolor{gray}{HTML}{c7c1bc}

\pgfkeys{
    /square/.is family, /square,
    default/.style = {color = {white}, border = {0.0}, fill = {full}},
    color/.store in = \squarecolor,
    border/.store in = \squareborder,
    fill/.store in = \squarefill
}

\NewDocumentCommand{\square}{o}{
    \pgfkeys{/square, default, #1}
    \begin{tikzpicture}[baseline={([yshift=-0.5em]current bounding box.center)}]
        \def\a{2em}

        \pgfmathparse{\squareborder > 0 ? "black" : "none"}
        \let\bordercolor\pgfmathresult
        \draw[fill=\squarecolor, draw=\bordercolor, line width=0.1em] (0,0) rectangle (\a,\a);
    \end{tikzpicture}
}

\pgfkeys{
    /diam/.is family, /diam,
    default/.style = {color = {white}, border = {0.0}, fill = {full}},
    color/.store in = \diamcolor,
    border/.store in = \diamborder,
    fill/.store in = \diamfill
}

\NewDocumentCommand{\diam}{o}{
    \pgfkeys{/diam, default, #1}
    % #1 = fill color, #2 = line width
    \begin{tikzpicture}[baseline={([yshift=-0.5em]current bounding box.center)}]
        \def\a{2.5em}
        \def\b{2em}

        \pgfmathparse{\diamborder > 0 ? "black" : "none"}
        \let\bordercolor\pgfmathresult
        \draw[fill=\diamcolor, draw=\bordercolor, line width=0.1em] (0,\b/2) -- (\a/2,0) -- (\a,\b/2) -- (\a/2,\b) -- cycle;
    \end{tikzpicture}
}

\pgfkeys{
    /olive/.is family, /olive,
    default/.style = {color = {white}, border = {0.0}, fill = {full}},
    color/.store in = \olivecolor,
    border/.store in = \oliveborder,
    fill/.store in = \olivefill
}

\NewDocumentCommand{\olive}{o}{
    \pgfkeys{/olive, default, #1}
    \begin{tikzpicture}[baseline={([yshift=-0.5em]current bounding box.center)}]
        \pgfmathsetmacro{\a}{2/1cm*1em}
        \pgfmathsetmacro{\b}{1.4/1cm*1em}
        \pgfmathparse{(\a*\a + \b*\b)/4/\b}
        \let\r\pgfmathresult
        \pgfmathparse{atan2(2*\r-\b, \a)}
        \let\x\pgfmathresult

        \pgfmathparse{\oliveborder > 0 ? "black" : "none"}
        \let\bordercolor\pgfmathresult
        \draw[fill=\olivecolor, draw=\bordercolor, line width=0.1em]
        (\a/2,\r-\a/2) arc[start angle=\x, end angle=180-\x, x radius=\r, y radius=\r]
        (\a/2,\r-\a/2) arc[start angle=-\x, end angle=-180+\x, x radius=\r, y radius=\r];
    \end{tikzpicture}
}

\pgfkeys{
    /tri/.is family, /tri,
    default/.style = {color = {white}, border = {0.0}, fill = {full}},
    color/.store in = \tricolor,
    border/.store in = \triborder,
    fill/.store in = \trifill
}

\NewDocumentCommand{\tri}{o}{
    \pgfkeys{/tri, default, #1}
    \begin{tikzpicture}[baseline={([yshift=-0.5em]current bounding box.center)}]
        \def\a{2em}
        \def\b{1.5em}

        \pgfmathparse{\triborder > 0 ? "black" : "none"}
        \let\bordercolor\pgfmathresult
        \draw[fill=\tricolor, draw=\bordercolor, line width=0.1em]
        (-\a/2, 0) -- (\a/2, \b/2) -- (\a/2, -\b/2) -- cycle;
    \end{tikzpicture}
}

\NewDocumentCommand{\pill}{m O{0.0}}{
    \begin{tikzpicture}[baseline={([yshift=-0.5em]current bounding box.center)}]
        \def\a{2em}
        \def\r{0.8em}

        \pgfmathparse{#2 > 0 ? "black" : "none"}
        \let\bordercolor\pgfmathresult
        \draw[fill=#1, draw=\bordercolor, line width=0.1em]
        (\r, 0) -- (\r+\a, 0) arc[start angle=-90, end angle=90, radius=\r] --
        (\r+\a, 2*\r) -- (\r, 2*\r) arc[start angle=90, end angle=270, radius=\r] -- cycle;
    \end{tikzpicture}
}

\NewDocumentCommand{\card}{m m m m}{
    % #1 = width, #2 = height, #3 = label
    \begin{tikzpicture}
        \def\barh{4ex}
        \def\r{0.5em}

        % background
        \draw[rounded corners=0.5em, draw=black, line width=0.1em ]
        (0, 0) rectangle(#1, #2);

        % top bar
        \draw[thick, fill=blue!10]
        (0, #2-\barh) -- (0, #2-\r) arc[start angle=180, end angle=90, radius=\r] --
        (\r, #2) -- (#1-\r, #2) arc[start angle=90, end angle=0, radius=\r] --
        (#1, #2-\r) -- (#1, #2-\barh) -- cycle;
        \node[anchor=west, font=\bfseries, xshift=0.5em, yshift=-0.2em] at (0, #2-\barh/2) {#3};

        \node[anchor=north, align=center, text width=#1-1em] at (#1/2, #2-\barh) {#4};
    \end{tikzpicture}
}